\documentclass[11pt,a4paper]{article}
\usepackage[utf8]{inputenc}
\usepackage[T1]{fontenc}
\usepackage{amsmath,amsfonts,amssymb}
\usepackage{graphicx}
\usepackage{hyperref}
\usepackage{listings}
\usepackage[dvipsnames]{xcolor}
\usepackage{booktabs}
\usepackage{float}
\usepackage{geometry}
\geometry{margin=1in}

\definecolor{zenblue}{RGB}{41,121,255}
\definecolor{zengreen}{RGB}{52,199,89}
\definecolor{codegray}{RGB}{245,245,245}

\hypersetup{colorlinks=true,linkcolor=zenblue,urlcolor=zenblue,citecolor=zenblue}

\lstset{
    backgroundcolor=\color{codegray},
    basicstyle=\ttfamily\small,
    breaklines=true,
    captionpos=b,
    frame=single,
    numbers=left,
    numberstyle=\tiny\color{gray}
}

\title{
    \vspace{-2cm}
    \Large \textbf{Zen AI Model Family} \\
    \vspace{0.5cm}
    \Huge \textbf{Zen-Director} \\
    \vspace{0.3cm}
    \large Single-GPU Unified Text- and Image-to-Video Generation \\
    \vspace{0.2cm}
    \large \normalsize A Packaging of Alibaba's Wan2.2-TI2V-5B (Apache-2.0) \\
    \vspace{0.5cm}
    \normalsize Technical Report
}

\author{
    Hanzo AI Research Team\thanks{research@hanzo.ai} \and
    Zoo Labs Foundation\thanks{foundation@zoo.ngo}
}

\date{2026}

\begin{document}

\maketitle

\begin{abstract}
\textbf{Zen-Director} is a permissively licensed redistribution and integration of
\textbf{Wan2.2-TI2V-5B}, the open-source unified text- and image-to-video generation model
released by the Alibaba Wan team (\texttt{Wan-AI/Wan2.2-TI2V-5B}) under the
\textbf{Apache-2.0 license}~\cite{wan22hf,wan22github}. Zen-Director is \emph{not} a
from-scratch model: it packages Wan2.2-TI2V-5B---a 5-billion-parameter \emph{dense} diffusion
transformer trained with flow matching, paired with the high-compression Wan2.2 3D causal
VAE---into the Zen model family with consistent tooling, weights distribution, and
serving integration. The defining property of Wan2.2-TI2V-5B is that it generates 720P video
at 24~fps on a \emph{single consumer-grade GPU} (e.g. an NVIDIA RTX~4090 with under 24~GB of
VRAM), generating a 5-second 720P clip in under nine minutes without specialized
optimization~\cite{wan22hf,wan22github}. A single unified framework natively serves both
text-to-video (T2V) and image-to-video (I2V). This report documents the model we redistribute,
its true architecture, its capabilities, and how it is integrated into the Zen stack. All
architectural facts and figures below are reported by the upstream Wan team; we add no
independent benchmarks.
\end{abstract}

\tableofcontents
\newpage

\section{Introduction}

Open-source video generation has reached the point where high-resolution, high-frame-rate
synthesis is feasible on a single consumer GPU. The Alibaba Wan team's \textbf{Wan2.2}
family~\cite{wan22github,wan21paper} is a leading example of this shift. Within that family,
\textbf{Wan2.2-TI2V-5B} is a compact 5-billion-parameter \emph{dense} model that unifies
text-to-video and image-to-video generation and runs within the memory budget of widely
available hardware. Crucially, both the code and the weights are released under the
\textbf{Apache-2.0 license}~\cite{wan22hf,wan22github}, which permits commercial redistribution
and modification with attribution.

\textbf{Zen-Director} is our packaging of Wan2.2-TI2V-5B. We make no claim to having designed,
trained, or fine-tuned a new architecture. Our contribution is integration: distributing the
upstream Apache-2.0 weights through the Zen model family, wiring the model into the Zen serving
and tooling stack, and providing a consistent interface alongside other Zen models. This
document therefore describes the \emph{real} Wan2.2-TI2V-5B architecture and capabilities so
that users of Zen-Director understand exactly what they are running and under what license.

\subsection{Provenance and License}

\begin{table}[H]
\centering
\begin{tabular}{ll}
\toprule
\textbf{Property} & \textbf{Value} \\
\midrule
Upstream model & Wan2.2-TI2V-5B \\
Upstream author & Alibaba Wan team (Wan-Video / Wan-AI) \\
Upstream repository & \texttt{Wan-AI/Wan2.2-TI2V-5B} (Hugging Face) \\
License & Apache-2.0 \\
Zen-Director role & Packaging / integration, not from-scratch training \\
Modality & Unified text-to-video (T2V) + image-to-video (I2V) \\
\bottomrule
\end{tabular}
\caption{Zen-Director provenance. Zen-Director redistributes the upstream model under its
original Apache-2.0 license with attribution to the Alibaba Wan team.}
\end{table}

\subsection{Model Overview}

\begin{table}[H]
\centering
\begin{tabular}{ll}
\toprule
\textbf{Property} & \textbf{Value (as reported upstream~\cite{wan22hf,wan22github})} \\
\midrule
Parameters & 5B (dense; not Mixture-of-Experts) \\
Architecture & Diffusion Transformer (DiT) + Wan2.2 3D causal VAE \\
Text encoder & umT5 (multilingual) \\
Training objective & Flow matching \\
VAE compression ($T\times H\times W$) & $4\times16\times16$ (overall rate $\approx 64$) \\
With patchification & $4\times32\times32$ \\
Output resolution & 720P ($1280\times704$ or $704\times1280$) \\
Frame rate & 24 fps \\
Clip length & up to 5 seconds \\
Text input & up to 512 tokens \\
Image input (I2V) & up to $1280\times704$ \\
GPU requirement & single consumer GPU, $<24$ GB VRAM (e.g. RTX 4090) \\
Reported speed & 5-second 720P clip in $<9$ minutes (single GPU, no special optimization) \\
\bottomrule
\end{tabular}
\caption{Wan2.2-TI2V-5B specifications as reported by the upstream Wan team. Zen-Director
redistributes these weights unchanged.}
\end{table}

\subsection{Key Capabilities}

\begin{itemize}
    \item \textbf{Unified text-to-video}: Generate a 720P, 24~fps clip (up to 5~seconds) from
    a text prompt of up to 512 tokens.
    \item \textbf{Unified image-to-video}: Animate a reference image (up to $1280\times704$)
    into a coherent video within the same single framework.
    \item \textbf{Single-GPU operation}: Run end-to-end on one consumer-grade GPU with under
    24~GB of VRAM, enabled by the high-compression 3D causal VAE.
    \item \textbf{Permissive licensing}: Apache-2.0 weights and code, suitable for commercial
    redistribution with attribution.
\end{itemize}

\section{Architecture}

Zen-Director runs the Wan2.2-TI2V-5B model unchanged. Its architecture follows the Wan design
lineage~\cite{wan21paper,wan22github}: a 3D causal variational autoencoder (VAE), a umT5 text
encoder, and a diffusion transformer (DiT), trained with flow matching. We describe each
component as documented upstream.

\subsection{5B Dense Diffusion Transformer}

The generative backbone is a \emph{dense} diffusion transformer with approximately 5 billion
parameters~\cite{wan22hf,wan22github}. Unlike the larger Wan2.2 A14B variants---which use a
Mixture-of-Experts design with a high-noise expert for early denoising and a low-noise expert
for later refinement---the TI2V-5B model is a single dense network with no expert routing. The
DiT operates on the compressed video latents produced by the VAE: a patchifying module forms
tokens from the latent, a stack of transformer blocks denoises them while attending to the text
embedding via cross-attention, and an unpatchifying module reconstructs the latent.

\subsection{High-Compression Wan2.2 3D Causal VAE}

The capability that lets a 720P 24~fps video model fit on a single 24~GB GPU is the
high-compression \textbf{Wan2.2-VAE}, a 3D causal autoencoder~\cite{wan22github,wan22hf}. It
achieves a spatio-temporal ($T\times H\times W$) compression ratio of $4\times16\times16$,
raising the overall information compression rate to roughly~64 while maintaining high-quality
reconstruction. Combined with the patchification applied inside the DiT, the effective spatial
reduction reaches $4\times32\times32$. The 3D \emph{causal} structure preserves temporal
causality across frames and substantially reduces the memory footprint of the latent sequence,
which is what makes single-GPU 720P generation practical.

\subsection{umT5 Text Encoder}

Text conditioning uses the umT5 multilingual text encoder~\cite{wan21paper}, which encodes the
input prompt (up to 512 tokens) into embeddings consumed by the DiT through cross-attention.
The Wan team selected umT5 for its strong multilingual (notably Chinese and English) encoding
and favorable training convergence.

\subsection{Flow-Matching Training Objective}

The diffusion transformer is trained with \textbf{flow matching}~\cite{wan21paper}. Given a
target latent $x_1$, a noise sample $x_0$, and a timestep $t \in [0,1]$, the training input is
the linear interpolation
\begin{equation}
    x_t = (1-t)\,x_0 + t\,x_1 ,
\end{equation}
and the model is trained to predict the velocity field
\begin{equation}
    v_t = x_1 - x_0 .
\end{equation}
At inference, samples are produced by integrating the learned velocity field, avoiding iterative
per-step velocity re-prediction schemes used by some alternative diffusion formulations.

\subsection{Unified T2V + I2V Framework}

A single model serves both modalities~\cite{wan22hf,wan22github}. For text-to-video the model is
conditioned on the umT5 text embedding alone; for image-to-video the reference image (up to
$1280\times704$) provides additional conditioning. Both paths share the same VAE, DiT, and
flow-matching sampler, so no separate model or checkpoint is required to switch between
text-driven and image-driven generation.

\section{Inference and Integration}

\subsection{Running Zen-Director}

Because Zen-Director is the upstream Wan2.2-TI2V-5B model, it is invoked through the standard
Wan2.2 generation interface. The example below illustrates the conceptual call (refer to the
upstream repository for the authoritative API and flags)~\cite{wan22github}.

\begin{lstlisting}[language=Python, caption=Invoking Zen-Director (Wan2.2-TI2V-5B)]
# Text-to-video: 720P, 24 fps, up to 5 seconds, single 24GB GPU
generate(
    task="ti2v-5B",
    size="1280*704",
    prompt="a lighthouse at dusk, waves breaking, slow push-in",
)

# Image-to-video: animate a reference image in the same framework
generate(
    task="ti2v-5B",
    size="1280*704",
    image="reference.png",
    prompt="gentle camera drift, wind moving through grass",
)
\end{lstlisting}

\subsection{Hardware Footprint}

The upstream Wan team reports that Wan2.2-TI2V-5B generates a 5-second 720P (24~fps) clip in
under nine minutes on a single consumer-grade GPU, without specialized optimization, within a
sub-24~GB VRAM budget such as an NVIDIA RTX~4090~\cite{wan22hf,wan22github}. This single-GPU
profile is the primary reason Zen-Director is suitable for self-hosted and on-premise
deployments where multi-GPU clusters are unavailable.

\subsection{Role in the Zen Stack}

Within the Zen model family, Zen-Director is the video-generation component. It consumes a text
prompt (and optionally a reference image) and produces video, complementing other Zen models
that handle language and image tasks. Integration consists of distributing the Apache-2.0
weights, exposing the unified T2V/I2V interface through Zen tooling, and providing consistent
serving alongside the rest of the family. No retraining or architectural modification is
performed.

\section{Reported Results}

The Wan team positions Wan2.2-TI2V-5B as one of the fastest 720P@24fps open models currently
available and reports favorable quality on its internal Wan-Bench~2.0 evaluation relative to
contemporary models~\cite{wan22github}. The public model card does not publish a specific
per-metric VBench table for the TI2V-5B variant, and \emph{we deliberately report no
independent benchmark numbers of our own}: Zen-Director redistributes the upstream weights
without modification, so any quantitative claims should be taken from, and attributed to, the
upstream Wan releases~\cite{wan22hf,wan22github,wan21paper}.

\section{Licensing and Attribution}

Zen-Director is governed by the \textbf{Apache-2.0 license} of the upstream
Wan2.2-TI2V-5B release~\cite{wan22hf,wan22github}. Redistribution under the Zen name retains
this license and includes attribution to the Alibaba Wan team as the originating authors of the
model, weights, and architecture. Apache-2.0 permits commercial use, modification, and
redistribution provided the license and notices are preserved. Users remain responsible for
ensuring that generated content complies with applicable laws and does not infringe third-party
rights.

\section{Related Work}

Wan2.2-TI2V-5B belongs to the broader Wan family of open large-scale video generative
models~\cite{wan21paper,wan22github}, which also includes the larger A14B Mixture-of-Experts
text-to-video and image-to-video models. The Wan line builds on the diffusion-transformer
approach to video synthesis, combining a 3D causal VAE for efficient latent representation, a
multilingual text encoder, and flow-matching training. Zen-Director's specific value is in
making the compact, single-GPU TI2V-5B configuration of this family conveniently available and
integrated within the Zen stack.

\section{Conclusion}

Zen-Director is an honest redistribution of Alibaba's Apache-2.0 Wan2.2-TI2V-5B model: a 5B
dense diffusion transformer with a high-compression 3D causal VAE that generates 720P, 24~fps
video from text or images on a single consumer GPU. It is not a from-scratch architecture and
introduces no new benchmarks; its contribution is packaging and integration. By attributing the
upstream Wan team and preserving the Apache-2.0 license, Zen-Director brings state-of-the-art,
permissively licensed, single-GPU video generation into the Zen model family with full
transparency about its provenance.

\begin{thebibliography}{10}
\bibitem{wan22hf} Alibaba Wan team, ``Wan2.2-TI2V-5B,'' Hugging Face model card, \texttt{Wan-AI/Wan2.2-TI2V-5B}. \url{https://huggingface.co/Wan-AI/Wan2.2-TI2V-5B}. License: Apache-2.0.
\bibitem{wan22github} Wan-Video, ``Wan2.2: Open and Advanced Large-Scale Video Generative Models,'' GitHub repository. \url{https://github.com/Wan-Video/Wan2.2}. License: Apache-2.0.
\bibitem{wan21paper} Wan team, ``Wan: Open and Advanced Large-Scale Video Generative Models,'' arXiv:2503.20314, 2025. \url{https://arxiv.org/abs/2503.20314}.
\end{thebibliography}

\end{document}
